\documentclass[11pt]{article}
\usepackage[margin=0.9in]{geometry}
\usepackage{amsmath,amssymb,amsthm,mathtools}
\usepackage[T1]{fontenc}
\usepackage[utf8]{inputenc}
\usepackage{enumitem}
\usepackage{booktabs,longtable,array,seqsplit}
\usepackage{hyperref}
\usepackage{xcolor}
\usepackage{microtype}
\hypersetup{colorlinks=true,linkcolor=blue!45!black,citecolor=blue!45!black,urlcolor=blue!45!black}
\setlength{\emergencystretch}{6em}

\newtheorem{theorem}{Theorem}[section]
\newtheorem{proposition}[theorem]{Proposition}
\newtheorem{lemma}[theorem]{Lemma}
\newtheorem{corollary}[theorem]{Corollary}
\newtheorem{claim}[theorem]{Claim}
\theoremstyle{definition}
\newtheorem{definition}[theorem]{Definition}
\newtheorem{assumption}[theorem]{Assumption}
\theoremstyle{remark}
\newtheorem{remark}[theorem]{Remark}

\newcommand{\R}{\mathbb{R}}
\newcommand{\C}{\mathbb{C}}
\newcommand{\N}{\mathbb{N}}
\newcommand{\Hsp}{\mathcal{H}}
\newcommand{\Fsp}{\mathcal{F}}
\newcommand{\Id}{\mathrm{Id}}
\newcommand{\tr}{\operatorname{tr}}
\newcommand{\diag}{\operatorname{diag}}
\newcommand{\rank}{\operatorname{rank}}
\newcommand{\spec}{\operatorname{spec}}
\newcommand{\spanop}{\operatorname{span}}
\newcommand{\dist}{\operatorname{dist}}
\newcommand{\range}{\operatorname{ran}}
\newcommand{\Ker}{\operatorname{ker}}
\newcommand{\Pinch}{\mathcal{C}}
\newcommand{\Off}{\widetilde{\mathcal{C}}}
\newcommand{\Pperp}{P^{\perp}}
\newcommand{\Qperp}{Q^{\perp}}
\newcommand{\op}{\mathrm{op}}
\newcommand{\F}{\mathrm{F}}
\newcommand{\HS}{\mathrm{HS}}
\newcommand{\ip}[2]{\left\langle #1,#2\right\rangle}
\newcommand{\norm}[1]{\left\lVert #1\right\rVert}
\newcommand{\abs}[1]{\left\lvert #1\right\rvert}
\newcommand{\code}[1]{\nolinkurl{#1}}

\newenvironment{argument}{\paragraph{Argument.}\begin{enumerate}[leftmargin=2em]}{\end{enumerate}}
\newenvironment{proofoutline}{\paragraph{Proof architecture.}\begin{enumerate}[leftmargin=2em]}{\end{enumerate}}
\newcommand{\sourceanchor}[1]{\par\smallskip\noindent\textbf{Source anchor.} #1\par\smallskip}
\newcommand{\sourcecomment}[1]{\begin{remark}#1\end{remark}}
\newenvironment{sourceguide}{\begin{description}[style=nextline,leftmargin=3.3cm,labelwidth=3.1cm]}{\end{description}}
\newenvironment{formalizationmap}{\begin{longtable}{@{}p{0.20\textwidth}p{0.31\textwidth}p{0.40\textwidth}@{}}\toprule Source item & Modern mathematical content & Repository status \\\midrule\endfirsthead\toprule Source item & Modern mathematical content & Repository status \\\midrule\endhead}{\bottomrule\end{longtable}}

\title{The sharp reciprocal-Fourier Sylvester estimate in Albeverio--Makarov--Motovilov (2001)}
\author{}
\date{Source-order distilled reconstruction from arXiv source math/0105142v3}
\begin{document}
\maketitle

\section*{Source map and scope}
\begin{sourceguide}
\item[Primary source] Sergio Albeverio, Konstantin A. Makarov, and Alexander K. Motovilov, \emph{Graph Subspaces and the Spectral Shift Function}, arXiv:math/0105142v3 (2001).
\item[Local evidence] Supplied compressed arXiv source \code{arXiv-math0105142v3.gz}.\\
Compressed SHA--256 \texttt{\seqsplit{40b4529a750165605c96856fd304dc17a8c57fb7619509f52f2289212a5a42d3}}.\\
Decompressed \code{ssflanl.tex} SHA--256 \texttt{\seqsplit{0ab8682761f4db2dfa0e3098bdd2887ae076d371a64816d808fa0a4dec7a2c86}}.
\item[Redistribution] The supplied source is not vendored. This document is a new mathematical reconstruction and formalization map.
\item[Coverage] Section ``Sylvester equation,'' especially Definition~\texttt{enorm}, Definition~\texttt{DefSolSyl}, Lemmas~\texttt{Krein} and \texttt{Kexp}, Theorem~\texttt{Kiexp}, equations \texttt{Tsemib}, \texttt{HeinzMod}, \texttt{fourie}, \texttt{ozenka}, \texttt{cpi2}, and Remark~\texttt{nad}. The later Riccati fixed-point use of the same Fourier extremal problem is included only where it confirms scaling.
\item[Formalization target] \code{ForMathlib/Analysis/InnerProductSpace/DavisKahanTheory/ReciprocalMultiplier.lean} and \code{Sylvester.lean}, especially the normalized finite reciprocal-Fourier interpolation seam and the sharp rectangular Ky Fan/UI-norm consequences.
\item[Attribution boundary] The paper is a bridge source for the sharp theorem, not the original source of every ingredient. It attributes the earlier spectral-subspace/Sylvester estimate to Bhatia--Davis--McIntosh (1983), the extremal Fourier norm $\pi/2$ to Sz.-Nagy (1953, with Strausz credited in the prose), and sharpness of the final constant to McEachin (1993).
\end{sourceguide}

The full paper concerns graph subspaces, Riccati equations, admissible operators, and spectral shift functions. The present reconstruction deliberately isolates the theorem family that is active in the finite Davis--Kahan formalization: solving a Sylvester equation between self-adjoint operators with disjoint spectra, representing the reciprocal spectral multiplier by Fourier characters, and obtaining the sharp $\pi/2$ bound in every symmetric operator ideal.

\section{The source setting}

\sourceanchor{Section ``Sylvester equation,'' opening paragraphs and Definition \texttt{enorm}.}

Let $\Hsp$ and $\Fsp$ be Hilbert spaces. The source writes the Sylvester equation as
\begin{equation}
  XA-CX=Y,
  \tag{AMM-Syl}\label{eq:amm-syl}
\end{equation}
where $A$ acts on $\Hsp$, $C$ acts on $\Fsp$, and $X,Y:\Hsp\to\Fsp$ are bounded. Both $A$ and $C$ may be unbounded in the paper, so the source distinguishes weak, strong, and operator solutions. The finite-dimensional formalization has bounded maps everywhere and therefore needs only the algebraic equation, but the source distinctions explain why its integral representation is initially stated in the weak operator topology.

A symmetric normed ideal $\mathcal S\subseteq\mathcal B(\Hsp,\Fsp)$ carries a norm $\lVert\!\lVert\!\lVert\cdot\rVert\!\rVert\!\rVert$ satisfying
\[
  \lVert\!\lVert\!\lVert KTH\rVert\!\rVert\!\rVert
  \le \norm{K}_{\op}\,
      \lVert\!\lVert\!\lVert T\rVert\!\rVert\!\rVert\,
      \norm{H}_{\op},
\]
and agrees with the operator norm on rank-one maps. The paper additionally assumes weak-operator lower semicontinuity for uniformly norm-bounded sequences. In finite dimensions this closure condition is automatic, while the first inequality is precisely the two-sided ideal property encoded by the repository's rectangular unitarily invariant norms.

\section{Three solution representations before the sharp theorem}

\subsection{Contour representation}

\sourceanchor{Lemma \texttt{Krein}, equation \texttt{ESylKrein}.}

If $A$ is closed, $C$ is bounded, and their spectra are disjoint, the source recalls the contour formula
\[
  X=\frac{1}{2\pi i}\int_\gamma
      (C-\zeta)^{-1}Y(A-\zeta)^{-1}\,d\zeta.
\]
The resulting ideal-norm estimate depends on the chosen contour length and resolvent bounds. This proves existence and uniqueness but does not expose the sharp universal function of spectral distance.

\subsection{Ordered half-plane representation}

\sourceanchor{Lemma \texttt{Kexp}, equation \texttt{XYSol1}.}

Under a one-sided accretivity hypothesis separated by $d>0$, Heinz's semigroup formula gives
\[
  X=\int_0^\infty e^{tC}Ye^{-tA}\,dt,
  \qquad
  \lVert\!\lVert\!\lVert X\rVert\!\rVert\!\rVert
  \le \frac1d\lVert\!\lVert\!\lVert Y\rVert\!\rVert\!\rVert.
\]
This is the infinite-dimensional analogue of the repository's sharp constant-one ordered and interval/exterior Sylvester theorems. It is stronger than the arbitrary-disjoint-spectrum estimate because the spectra are separated by an orientation, not merely by absolute distance.

\subsection{Why arbitrary disjoint spectra are different}

If the spectra interlace or lie on both sides of one another, no single positive semigroup parameter resolves every denominator with the same sign. In joint eigen-coordinates the equation is
\[
  (\alpha_i-\beta_j)X_{ij}=Y_{ij}.
\]
The solution is therefore the Schur multiplier
\[
  X_{ij}=\frac{1}{\alpha_i-\beta_j}Y_{ij},
\]
and the problem becomes one of representing the reciprocal function on the separated difference set by averages of characters.

\section{The sharp separated-spectrum theorem}

\sourceanchor{Theorem \texttt{Kiexp}, equations \texttt{Tsemib}, \texttt{HeinzMod}, \texttt{fourie}, \texttt{ozenka}, \texttt{cpi2}, and \texttt{ozenkaR}.}

Assume $A=A^*$ and $C=C^*$ and set
\[
  d:=\dist\bigl(\spec(A),\spec(C)\bigr)>0.
\]
The paper states that the unique weak solution of \eqref{eq:amm-syl} has a Fourier-orbit representation
\begin{equation}
  X=\int_{-\infty}^{\infty}
      e^{itC}Ye^{-itA}f_d(t)\,dt.
  \tag{AMM-Fourier}
\end{equation}
Here $f_d\in L^1(\R)$ is chosen so that its Fourier transform agrees with the reciprocal function on all spectral differences. With a consistent scaling convention this condition is
\begin{equation}
  \widehat f_d(x)=\int_{\R}e^{-itx}f_d(t)\,dt=\frac1x
  \qquad\text{whenever }\abs{x}\ge d.
  \tag{AMM-Reciprocal}
\end{equation}
Indeed, taking the $(i,j)$ matrix coefficient of \eqref{AMM-Fourier} multiplies $Y_{ij}$ by $\widehat f_d(\alpha_i-\beta_j)$.

Each left and right exponential is unitary. Thus every symmetric ideal norm satisfies
\[
  \lVert\!\lVert\!\lVert X\rVert\!\rVert\!\rVert
  \le \norm{f_d}_{L^1}\,
      \lVert\!\lVert\!\lVert Y\rVert\!\rVert\!\rVert.
\]
The extremal reciprocal-Fourier problem gives kernels with mass arbitrarily close to
\[
  \frac{\pi}{2d}.
\]
Consequently
\begin{equation}
  \lVert\!\lVert\!\lVert X\rVert\!\rVert\!\rVert
  \le \frac{\pi}{2d}
      \lVert\!\lVert\!\lVert Y\rVert\!\rVert\!\rVert.
  \tag{AMM-Sharp}
\end{equation}
The operator norm is a special case.

\subsection{Verification of the Sylvester identity}

For bounded finite-dimensional operators, the formal verification is transparent. Define
\[
  \Phi(Y):=\int_{\R}e^{itC}Ye^{-itA}f_d(t)\,dt.
\]
In eigenbases $Ce_i=\alpha_i e_i$ and $Au_j=\beta_j u_j$,
\[
  \ip{e_i}{\Phi(Y)u_j}
  =\widehat f_d(\beta_j-\alpha_i)\ip{e_i}{Yu_j},
\]
with the sign determined by the Fourier convention. Choosing the convention to make this multiplier $(\beta_j-\alpha_i)^{-1}$ yields
\[
  \Phi(Y)A-C\Phi(Y)=Y.
\]
The source formula and the repository formula use opposite presentations of the Sylvester commutator in places; the formalization should normalize the sign once at the matrix-coefficient boundary rather than scatter sign repairs through the orbit argument.

\section{The extremal constant and its provenance}

\sourceanchor{Remark \texttt{nad}.}

The source separates three historical claims.

\begin{enumerate}[leftmargin=2em]
\item Bhatia--Davis--McIntosh (1983) proved the separated-spectrum theorem with bounds
\[
  \sqrt{\frac32}\le c\le 2
\]
for the best universal constant and reduced the upper estimate to the Fourier extremal quantity
\[
  \inf\left\{\norm{f}_{L^1(\R)}:
    \widehat f(x)=\frac1x\text{ for }\abs{x}\ge1\right\}.
\]
\item The source attributes the identity
\begin{equation}
  \inf\left\{\norm{f}_{L^1(\R)}:
    \widehat f(x)=\frac1x\text{ for }\abs{x}\ge1\right\}
  =\frac\pi2
  \tag{AMM-Favard}
\end{equation}
to the Bohr-inequality work of B.~Sz.-Nagy, with A.~Strausz credited in the prose.
\item McEachin (1993) proved that the resulting $\pi/2$ constant in the subspace perturbation bound is sharp.
\end{enumerate}

Therefore this paper is an excellent theorem-level bridge but must not be cited as if it supplied the original proof of \eqref{AMM-Favard} or the original sharpness construction.

\section{Scaling audit and a source-level inconsistency}

The displayed condition \texttt{fourie} in Theorem \texttt{Kiexp} says that
\[
  \widehat f_d(x)=1/x\quad\text{for }\abs{x}\ge 1/d.
\]
Taken literally together with $d=\dist(\spec(A),\spec(C))$, this does not produce the stated $1/d$ norm scaling.

Later, in the Riccati fixed-point proof, the source starts from a unit-gap kernel $f$ satisfying
\[
  \widehat f(s)=1/s\quad(\abs{s}\ge1)
\]
and defines
\[
  f_d(t):=f(dt).
\]
Then
\[
  \widehat f_d(x)=\frac1d\widehat f(x/d)=\frac1x
  \quad\text{when }\abs{x}\ge d,
  \qquad
  \norm{f_d}_{L^1}=\frac1d\norm{f}_{L^1}.
\]
This is consistent with \eqref{AMM-Sharp}. The reconstruction therefore treats $\abs{x}\ge d$ as the intended threshold and records the displayed $1/d$ threshold as a source-level scaling slip, not as a mathematical hypothesis to copy into Lean.

\section{Finite exact interpolation versus operator integration}

The paper uses an $L^1$ Fourier kernel and a weak operator integral. The finite formalization can avoid building a general operator-valued integration layer.

Let $\alpha:I\to\R$ and $\beta:J\to\R$ be the finite eigenvalue lists, with
\[
  \delta\le\abs{\alpha_i-\beta_j}.
\]
It is enough, for every $\varepsilon>0$, to construct finitely many coefficients $a_r\in\C$ and frequencies $t_r\in\R$ such that
\begin{align*}
  \delta
  &= (\alpha_i-\beta_j)
     \sum_r a_r e^{it_r(\alpha_i-\beta_j)},
     &&\text{for every }(i,j),\\
  \sum_r\abs{a_r}
  &\le \frac\pi2+\varepsilon.
\end{align*}
This is the repository predicate
\code{HasFiniteReciprocalFourierInterpolation}.

For complex Hilbert spaces, each character is realized by diagonal unitaries:
\[
  U_r e_i=e^{it_r\alpha_i}e_i,
  \qquad
  V_r u_j=e^{-it_r\beta_j}u_j.
\]
Their two-sided action multiplies the $(i,j)$ matrix unit by
$e^{it_r(\alpha_i-\beta_j)}$. The scalar interpolation therefore becomes one common finite orbit identity, from which every Ky Fan prefix and every rectangular unitarily invariant norm estimate follow.

The exact finite interpolation is stronger than pointwise approximation. Starting from an $L^1$ kernel, one still needs an exact finite moment-compression argument on the finite difference set, or an error-correction construction whose extra mass tends to zero. The source does not supply that finite reduction because its weak integral is already an acceptable endpoint.

\section{Why the real case is a separate formal obligation}

Complex diagonal phases are not diagonal real orthogonal maps unless the phases are signs. The complex proof therefore does not directly instantiate the repository's theorem quantified over an arbitrary real-or-complex scalar field.

A faithful real route should pass through one of the following equivalent mechanisms:

\begin{enumerate}[leftmargin=2em]
\item complexify the real spaces and prove that singular values and Ky Fan sums of the real operator agree with those of its complexification;
\item realize each complex phase by a two-dimensional real rotation on a doubled space, then compress or identify the original singular values;
\item prove the real Ky Fan inequality by a field-independent dilation theorem and recover any real orbit-convexity statement only afterward.
\end{enumerate}

Requiring an exact finite orbit certificate on the original real spaces before proving the Ky Fan estimate may be stronger than the source argument and stronger than the application requires.

\section{Formalization map}

\begin{formalizationmap}
Theorem \texttt{Kiexp}, spectral distance $d>0$ & General separated-spectrum Sylvester equation & \code{SpectraSeparated} and \code{kyFan_sylvester_le_of_spectralDistance} in \code{Sylvester.lean}; the generic sharp chain still depends on the open harmonic root. \\
Equation \texttt{fourie} & Reciprocal Fourier multiplier on the spectral difference set & \code{HasFiniteReciprocalFourierInterpolation} in \code{ReciprocalMultiplier.lean}; the formal statement uses a finite exact character sum. \\
Equation \texttt{HeinzMod} & Two-sided unitary orbit average & \code{basisDiagonalUnitary}, \code{unitaryOrbitAction_basisMatrixUnit}, and \code{complexUnitaryOrbitAction_basisMatrixUnit_exp_sub}. \\
Equation \texttt{ozenka} & Symmetric-ideal norm domination & Finite orbit certificate $\Rightarrow$ simultaneous rectangular Ky Fan bounds $\Rightarrow$ \code{uiNorm_sylvester_le_of_spectralDistance}. \\
Equation \texttt{cpi2} and Remark \texttt{nad} & Sharp constant as an infimum of Fourier masses & \code{kyFan_reciprocalMultiplier_le_complex_of_approximateFourierInterpolation}; it correctly needs masses $\pi/2+\varepsilon$, not an attaining kernel. \\
Unit-gap to gap-$d$ scaling & $f_d(t)=f(dt)$ and $\norm{f_d}_1=d^{-1}\norm f_1$ & \code{hasFiniteReciprocalFourierInterpolation_of_normalized}. \\
Symmetric ideal property & Invariance under left/right unitaries and ideal contraction & \code{RectangularUnitarilyInvariantNorm} and rectangular Fan dominance. \\
Real scalar field & Not supplied by complex diagonal phases & Still requires singular-value-preserving complexification or doubled-space rotation descent. \\
\end{formalizationmap}

\section{Current proof obligations exposed by the source}

The source makes the remaining work unusually precise.

\begin{enumerate}[leftmargin=2em]
\item Prove the normalized scalar theorem: for every finite separated difference set and every $\varepsilon>0$, construct an exact finite reciprocal-Fourier interpolation of mass at most $\pi/2+\varepsilon$.
\item Derive that construction from the Sz.-Nagy--Strausz extremal theorem or a modern equivalent, while preserving exactness on the finite set.
\item Transfer the sharp complex Ky Fan result to real Hilbert spaces without a constant loss.
\item Replace the field-uniform harmonic placeholder only after both field cases are connected to the generic theorem.
\end{enumerate}

No additional barycentric, residual, or tangent wrapper is needed before these obligations are discharged; the downstream finite operator theory is already in place.

\section*{Source-faithfulness and limitations}

This reconstruction follows the source order inside its Sylvester section and preserves the paper's distinction between ordered constant-one estimates and arbitrary-disjoint-spectrum $\pi/2$ estimates. It does not reconstruct the paper's main spectral-shift theorem, its full Stieltjes operator-integral machinery, or its graph-subspace/Riccati development except where the later fixed-point proof resolves the Fourier scaling convention.

The note also keeps attribution explicit. Albeverio--Makarov--Motovilov state and use the sharp theorem; Bhatia--Davis--McIntosh provide the earlier perturbation/Sylvester reduction; Sz.-Nagy--Strausz supply the extremal reciprocal-Fourier norm; McEachin supplies sharpness. A future reconstruction of Bhatia--Davis--McIntosh (1983) remains necessary to document the original operator-theoretic proof in its own source order.

\end{document}
